\section{Experimental Plan}
\label{sec:experimental-plan}

% =========================================================================
% Stub outline. Target length for first draft: ~2 pages.
% This section pre-registers the depth-limit ablation. It re-uses the
% framework paper's benchmarks and base models rather than introducing
% new ones, because the central claim is about depth-limit's effect on
% an otherwise-fixed system.
% =========================================================================

\subsection{Setup Inherited from the Framework Paper}
\label{sec:exp-inherited}

The benchmark, base-model, and verifier choices are inherited from our
framework paper without modification. We restate them briefly here for
self-containment.

\paragraph{Benchmarks.}
GSM8K \citep{cobbe2021gsm8k} and MATH \citep{hendrycks2021math} are the
primary evaluation targets. Both expose a strict-format terminal
verifier (regex-extracted boxed answer, SymPy normalization, equality
check). Competition mathematics (AIME) is a stretch target. Formal
proof (Lean via LeanDojo \citep{yang2023leandojo}) is deferred.

\paragraph{Base models.}
Qwen2.5-1.5B-Base \citep{yang2024qwen25} as the primary substrate, with
Llama-3.2-1B-Base \citep{dubey2024llama3} and SmolLM2-1.7B
\citep{allal2024smollm2} as secondary. We use base (non-instruct,
non-reasoning) variants so that any measured improvement is attributable
to the deliberation loop, not to post-training that already includes
chain-of-thought capability.

\paragraph{Architecture.}
The shared frozen LLM, adapter MLP, multi-head $Q$, policy head, and
verb-instantiator are specified by the framework paper. No
modifications.

% TODO: explicit reference to framework paper sections (\S4.4 architecture,
% \S4.3 action space, \S4.2 state schema) once the framework paper is
% citable. For now we reference by description.

\subsection{The Depth-Limit Ablation}
\label{sec:exp-ablation}

The headline ablation varies the per-hypothesis tree depth limit
$\depthlim$ across three settings, holding all other hyperparameters
fixed. Each setting is trained from scratch using the framework paper's
expert-iteration loop, with all other configuration identical to the
framework paper's $B4$ baseline (high simulation budget at inference).

\begin{table}[t]
\centering
\caption{Pre-registered depth-limit ablation grid. All three rows share
the framework paper's setup; only $\depthlim$ varies.}
\label{tab:depth-ablation}
\small
\begin{tabular}{@{}llp{0.55\linewidth}@{}}
\toprule
ID & $\depthlim$ & Description \\
\midrule
R0 & $0$       & No recursion. Equivalent to the framework paper's phase 1
                behavior. \textsc{Decompose} and \textsc{CheckCase} are
                permitted only at the top level; the
                \texttt{parent\_hypothesis} and analogous fields are
                never set. \\
R3 & $3$       & Productive-regime hypothesis. Bounded recursion: up to
                three nested levels of decomposition. Default value
                proposed in the framework paper. \\
R$\infty$ & $\infty$ & Unbounded recursion. No structural prior on
                tree depth. The recursive verbs remain available as
                long as the global depth budget $d_{\max}$ is not
                exhausted. \\
\bottomrule
\end{tabular}
\end{table}

\paragraph{Pre-registered hypotheses.}
We commit before training to the following three primary hypotheses
and one mechanism sub-hypothesis. Effect-size thresholds are stated
in pass@1 percentage points and are pre-registered together with the
statistical analysis plan in Section~\ref{sec:exp-stats}.
\begin{enumerate}[leftmargin=*]
  \item H1 (\emph{expressivity}): on MATH levels 1 through 3 (the
        benchmark slice with the most decomposable problem structure),
        R3 achieves higher pass@1 than R0 by at least one percentage
        point at every FLOP-matched evaluation point, averaged across
        at least three training seeds. Falsified if R0 matches R3
        (within one percentage point) at any of the
        FLOP-matched evaluation points.
  \item H2 (\emph{rumination}): R$\infty$ achieves lower pass@1 than
        R3 by at least one percentage point at every FLOP-matched
        evaluation point, averaged across at least three training
        seeds. Falsified if R$\infty$ matches or exceeds R3 (within
        one percentage point) at any of the FLOP-matched evaluation
        points.
  \item H2a (\emph{mechanism}): conditional on H2 being supported (R3
        outperforming R$\infty$ by at least the threshold), the MCTS
        budget allocation in R$\infty$ shows at least 25\% of search
        visits at hypothesis-tree depth $\geq 4$, and the
        verifier-call frequency at hypothesis-tree depth $\geq 4$ in
        R$\infty$ is below $0.05$ calls per trajectory. Falsified if
        R$\infty$'s underperformance occurs without this depth-tail
        signature in the search-tree statistics, in which case H2's
        mechanism attribution is unsupported even if its outcome
        inequality holds.
  \item H3 (\emph{depth distribution}): in successful R3 trajectories,
        the modal observed hypothesis-tree depth at
        \textsc{CommitAnswer} time is $1$ or $2$, with depth $3$ used
        only on the harder problem tail (operationalized as MATH
        problems in the top quartile by reported difficulty rating).
        Falsified if depth $0$ or depth $3$ dominate the modal
        distribution, indicating that the depth-limit mechanism is
        not exercising the intended regime.
\end{enumerate}

H1 and H2 jointly test the structural-prior story: the productive
regime should lie strictly inside the depth-zero and depth-unbounded
endpoints of the ablation grid. H2a separates outcome from mechanism
so that H2's causal attribution to budget dilution is itself
falsifiable. H3 tests whether the policy is actually using recursion
when it is available, as opposed to ignoring the new expressivity and
behaving like R0 with extra optionality.

\subsection{Secondary Measurements}
\label{sec:exp-secondary}

\paragraph{Trajectory composition.}
For each setting, we record the empirical distribution over tree depths
at decision-node level and at terminal-state level. This provides the
data for H3 and lets us see whether R$\infty$ exhibits long-tail
unproductive trajectories (the rumination signature: deep trees with
many \textsc{FormHypothesis} and \textsc{Decompose} calls and few
\textsc{InvokeVerifier} resolutions).

\paragraph{Verifier-call frequency by depth.}
We measure how often \textsc{InvokeVerifier} fires as a function of the
depth of its target. Productive recursion should produce a non-trivial
fraction of verifier calls at depth $\geq 1$; if all calls fire on
depth-$0$ entities, recursion is decorative.

\paragraph{MCTS budget allocation.}
For matched inference budget $B$, we record the depth distribution of
visited nodes in the search tree. Comparing R3 and R$\infty$ should
expose whether the unbounded-depth setting is concentrating budget along
unproductive deep paths (visible as a long tail in the depth
distribution at fixed $B$).

\subsection{Threats to Validity}
\label{sec:exp-threats}

The ablation grid inherits the framework paper's threats to validity
(compute accounting leaks, verifier coverage gaps, benchmark
contamination, seed variance) and adds four depth-specific threats
that the headline grid is not by itself equipped to rule out.

\paragraph{The depth limit might be doing nothing.}
The headline R3-versus-R0 contrast assumes the trained policy
actually exercises the recursive verbs that R3 unlocks. If the policy
under R3 never selects \textsc{Decompose}, \textsc{CheckCase}, or
recursive \textsc{FormHypothesis}, then R0 and R3 are indistinguishable
behaviorally and any pass@1 difference reflects training noise rather
than the depth prior. Hypothesis H3 (the depth-distribution
hypothesis) is the safeguard: if the modal observed tree depth at
\textsc{CommitAnswer} time in R3 is zero, the depth prior is not the
active mechanism and the ablation result is not informative about the
structural-recursion claim. We report the depth-at-commit distribution
alongside the headline pass@1 numbers so this failure mode is visible
even if H1 numerically passes.

\paragraph{R$\infty$ might fail for the wrong reason.}
Hypothesis H2 attributes R$\infty$'s expected underperformance to
simulation-budget dilution along operationally unproductive paths.
A competing explanation is training instability: the unbounded state
space inflates the parameter regime $\Qpsi$ must cover, which can
manifest as training-loss spikes, target-network drift, or gradient
pathologies independent of the rumination mechanism. We mitigate by
reporting per-iteration training diagnostics (loss curves,
target-network update magnitudes, gradient norms) for all three
settings, so that R$\infty$ underperformance attributable to training
instability is distinguishable from underperformance attributable to
rumination. Sub-hypothesis H2a (Section~\ref{sec:exp-ablation})
operationalizes this distinction: if R$\infty$'s budget allocation
does \emph{not} show the depth-tail signature we predict, the
mechanism claim does not hold even when the pass@1 inequality does.

\paragraph{A single $\depthlim$ value may not be the productive optimum.}
The headline R3 setting uses $\depthlim = 3$ on the basis of the
companion paper's qualitative argument
(Section~\ref{sec:method-depthlimit}), but the productive regime
might in fact lie at $\depthlim = 2$ or $\depthlim = 4$. A
three-setting grid (R0, R3, R$\infty$) cannot distinguish these
possibilities, and a single-point R3 result that disconfirms H1 may
indicate either ``no productive recursion regime exists'' or
``$\depthlim = 3$ is the wrong sample.'' We pre-register a secondary
fine-grained sweep over $\depthlim \in \{1, 2, 3, 4, 5\}$ on a
held-out 200-problem subset of MATH levels 1 through 3, run after the
headline three-setting result is finalized. The sweep is descriptive,
not inferential: we report the depth-versus-pass@1 curve without
additional hypothesis tests. Its purpose is to localize the
productive regime if the headline ablation supports H1, and to rule
out the ``wrong depth value'' explanation if the headline ablation
disconfirms it.

\paragraph{Verifier strength may dominate the structural prior.}
Productivity (Definition~\ref{def:productive}) is defined relative to
the verifier. If the symbolic verifier is too weak to verify the
intermediate claims that productive decomposition produces, no
depth-limit value will surface the productive regime. We mitigate by
reporting per-action verifier-call success rates by depth, so that a
verifier-coverage gap at depth $\geq 1$ becomes visible. A null
result on H1 accompanied by low verifier-call success rates at
depths $\geq 1$ does not falsify the productive-recursion claim; it
falsifies the joint hypothesis that the present verifier is strong
enough to expose it.

\subsection{Reporting Discipline}
\label{sec:exp-reporting}

Following the framework paper's pre-registration discipline:
\begin{itemize}[leftmargin=*]
  \item All three settings are trained to a fixed total compute budget,
        not to a fixed accuracy ceiling. Training compute is a
        controlled variable.
  \item All three settings use identical random seeds for the policy
        and $Q$-head initialization. Replay buffer samples are
        independent.
  \item Mean and standard deviation across at least three training
        seeds are reported for the headline pass@1 numbers.
  \item Compute accounting reports total inference FLOPs, total LLM
        forward passes, and wall-clock seconds at fixed hardware, as
        defined in the framework paper.
\end{itemize}

\subsection{Statistical Analysis Plan}
\label{sec:exp-stats}

Pre-registration without a statistical analysis plan reduces to a
sequence of qualitative claims that admit post-hoc selection of the
favorable comparison. We commit before training to the following
analysis plan, mirroring the framework paper's plan with the depth
ablation's specific structure.

\paragraph{Effect-size thresholds.}
Hypotheses H1 and H2 are pre-registered with a one-percentage-point
pass@1 effect-size threshold. ``R3 outperforms R0'' (H1) is
operationalized as $\overline{\Delta}_{\text{R3,R0}} \geq 0.01$ at
every FLOP-matched evaluation point, where $\overline{\Delta}$ is the
mean per-problem accuracy difference across at least three training
seeds; symmetric for H2 in the opposite direction. The one-point
threshold matches the framework paper's effect-size convention and is
conservative relative to typical MATH pass@1 variation across
training runs.

\paragraph{Inference compute budgets.}
Each setting is evaluated at $K = 3$ FLOP-matched inference compute
budgets per benchmark, corresponding to MCTS simulation budgets
$B \in \{100, 300, 1000\}$ for the search-mode condition. Fast-mode
($B = 0$) is reported as a fourth point for context but is not part
of the H1 or H2 inferential test (the depth limit only takes effect
when MCTS expands recursive verbs at the root, so the $B = 0$ rows
collapse the three settings to the same trained $\Qpsi$ acted on
greedily). The H1 and H2 inferential tests cover the three $B > 0$
points on the single benchmark (MATH levels 1 through 3), for
$3 \times 1 = 3$ paired tests per hypothesis.

\paragraph{Inference procedure.}
For each FLOP-matched evaluation point we compute a paired bootstrap
on the per-problem accuracy difference $\Delta_i$ with $R = 10{,}000$
resamples. The $95\%$ confidence interval is the $2.5$th and $97.5$th
percentile of the resampled $\overline{\Delta}$ distribution. The
null hypothesis at each evaluation point is
$\overline{\Delta} = 0$; rejection means the CI excludes zero in the
hypothesized direction.

\paragraph{Multiple-comparison correction.}
H1 and H2 each span three FLOP-matched evaluation points, for six
primary inferential tests. We apply Holm-Bonferroni correction at
family-wise error rate $\alpha = 0.05$ across the six tests jointly.
H1 (resp.\ H2) is supported only if all three of its evaluation
points reject in the predicted direction after correction. H2a is a
conditional mechanism test reported alongside H2; we do not subject
H2a to Holm-Bonferroni correction because its purpose is to qualify
the interpretation of H2 rather than to add an independent claim.
H3 is a descriptive depth-distribution claim and is reported as
modal-depth percentages with bootstrap CIs rather than as a
hypothesis test; it has no Holm-Bonferroni budget.

\paragraph{Seed variance.}
All three settings are trained from at least three independent random
seeds. We report mean and standard deviation across seeds for every
pass@1 point. The per-problem accuracy used in the bootstrap is the
mean across seeds; the confidence interval therefore reflects
problem-level variance rather than seed variance, and seed variance
is reported alongside as a transparency check on training stability
(also a safeguard for the training-instability threat in
Section~\ref{sec:exp-threats}).

\paragraph{Pre-registration scope.}
The decisions pre-registered here (effect-size thresholds, FLOP-matched
evaluation points, paired-bootstrap procedure, Holm-Bonferroni
correction, seed count) cannot be revised post-hoc. The fine-grained
depth sweep over $\depthlim \in \{1, 2, 3, 4, 5\}$
(Section~\ref{sec:exp-threats}) is descriptive rather than
inferential; its results may suggest follow-up experiments but cannot
be retroactively used to substitute for the headline R3 result.
Failure of H1 or H2 under the pre-registered plan is reported as a
falsification of the corresponding claim, not as a search for a
modified plan under which the claim would have been supported.
